\documentclass[a4paper,12pt]{article}
\usepackage[T2A]{fontenc}
\usepackage[utf8]{inputenc}
\usepackage[russian]{babel}
\usepackage{amsmath,amssymb,mathtools}
\usepackage{geometry}
\usepackage{microtype}
\usepackage{array,booktabs,tabularx}
\usepackage{xcolor}
\usepackage{tikz}
\usepackage{hyperref}
\usepackage{enumitem}
\usepackage{fancyhdr}
\usepackage{setspace}
\geometry{left=22mm,right=22mm,top=23mm,bottom=25mm}
\definecolor{accent}{HTML}{087F83}
\definecolor{darkaccent}{HTML}{075C61}
\definecolor{textgray}{HTML}{444444}
\definecolor{err}{HTML}{AD3535}
\hypersetup{colorlinks=true,linkcolor=darkaccent,pdfauthor={Лёвин Артём Александрович},pdftitle={Ревью домашней работы — Анна Баннова — 15.11.2026}}
\setlength{\parindent}{0pt}
\setlength{\parskip}{7pt}
\setlength{\headheight}{14pt}
\setstretch{1.14}
\pagestyle{fancy}
\fancyhf{}
\fancyhead[L]{\small\color{textgray} Ревью домашней работы}
\fancyhead[R]{\small\color{textgray} Анна Баннова}
\fancyfoot[C]{\small\color{textgray}\thepage}
\renewcommand{\headrulewidth}{0.3pt}
\renewcommand{\footrulewidth}{0pt}
\newcommand{\subtopic}[1]{\par\medskip\noindent{\bfseries\color{darkaccent}#1}\par\vspace{-2pt}}
\begin{document}
\thispagestyle{empty}
\begin{center}
\vspace*{24mm}
{\large\color{darkaccent} МАТЕМАТИКА}\par
\vspace{24mm}
{\LARGE Ревью домашней работы}\par
\vspace{15mm}
{\Large\color{accent} Движение по круговой трассе}\par
\vspace{7mm}
{\normalsize Задачи на движение}\par
\vspace{47mm}
{\large Анна Баннова}\par
\vspace{12mm}
{\normalsize 15 ноября 2026 г.}\par
\vspace{19mm}
{\small\color{textgray} Лёвин Артём Александрович}\par
{\small\color{textgray} эксклюзивно для Анны Банновой}\par
\vfill
\begin{tikzpicture}[x=1cm,y=1cm]
 \draw[accent,line width=0.65pt] (0,0) .. controls (1.1,0.55) and (2.4,-0.5) .. (3.7,0)
   .. controls (4.8,0.45) and (5.8,-0.42) .. (7.0,0);
 \fill[accent] (3.7,0) circle (1.1pt);
\end{tikzpicture}
\vspace*{8mm}
\end{center}
\clearpage
\section*{Сводная таблица}
\addcontentsline{toc}{section}{Сводная таблица}
\smallskip
В таблице приведено задание, в котором обнаружена ошибка. Нажатие на номер переносит к подробному разбору.\par
\vspace{6mm}
\begingroup
\renewcommand{\arraystretch}{1.45}
\small
\begin{tabularx}{\linewidth}{@{}>{\centering\arraybackslash}p{13mm} >{\raggedright\arraybackslash}p{28mm} >{\raggedright\arraybackslash}X >{\centering\arraybackslash}p{22mm} >{\centering\arraybackslash}p{26mm}@{}}
\toprule
\textbf{№} & \textbf{Тема} & \textbf{Суть ошибки} & \textbf{Ваш ответ} & \textbf{Правильный ответ}\\
\midrule
\hyperref[task:2]{\textbf{2}} & Движение по кругу & Перепутан порядок вычитания расстояний при обгоне: от пути первого автомобиля вычтен путь второго. & $66$ км/ч & $102$ км/ч\\
\bottomrule
\end{tabularx}\par
\endgroup
\vspace{9mm}
\textcolor{darkaccent}{\textbf{На что обратить внимание.}}
Когда второй участник догоняет первого и опережает его на полный круг, именно \emph{второй} проходит на одну длину трассы больше. Этот смысл нужно зафиксировать до составления уравнения.
\clearpage
\thispagestyle{plain}
\vspace*{0.37\textheight}
\begin{center}
{\LARGE\color{darkaccent} Разбор ошибок}\par
\vspace{5mm}
{\normalsize\color{textgray} Движение по круговой трассе}
\end{center}
\clearpage
\phantomsection\label{task:2}
\section*{Задание № 2. Поиск скорости автомобиля}
\subtopic{Текст задания}
Из одной точки круговой трассы длиной $12$ км одновременно выехали в одном направлении два автомобиля. Первый движется с постоянной скоростью $84$ км/ч. Через $40$ минут после старта он впервые был обогнан вторым автомобилем ровно на один круг. Найдите постоянную скорость второго автомобиля. Ответ дайте в километрах в час.

\subtopic{Ваше решение}
Вы верно перевели время в часы и выразили расстояния, которые успели пройти автомобили. Обозначив скорость второго автомобиля через $x$, Вы получили:
\[
 t=\frac{40}{60}=\frac23\text{ ч},\qquad
 S_1=84\cdot\frac23=56\text{ км},\qquad
 S_2=x\cdot\frac23=\frac{2x}{3}\text{ км}.
\]
Далее в тетради записано:
\[
 56-\frac{2x}{3}=12,
 \qquad \frac{2x}{3}=44,
 \qquad x=66.
\]

\subtopic{Ваш ответ}
$66$ км/ч.

\subtopic{Ошибка}
\textcolor{err}{\textbf{Ошибка:}} неверно выбран порядок вычитания пройденных расстояний.

\subtopic{Где именно возникает ошибка}
\textbf{Последний правильный шаг:} $S_1=56$ км и $S_2=\dfrac{2x}{3}$ км.

\textbf{Первый неправильный шаг:}
\[
 \textcolor{err}{56-\frac{2x}{3}=12}.
\]
Из этого равенства следует, что \emph{первый} автомобиль прошёл на $12$ км больше второго. Однако по условию \emph{второй} автомобиль обогнал первый на целый круг.

\subtopic{Почему это неверно}
Длина трассы равна $12$ км. При старте из одной точки и движении в одну сторону обгон на один полный круг означает, что более быстрый участник прошёл на $12$ км \emph{больше}, чем другой. Следовательно, к моменту обгона путь второго автомобиля превосходит путь первого:
\[
 S_2-S_1=12.
\]
Разность $S_1-S_2=12$ описывала бы противоположную ситуацию: первый автомобиль обогнал второй на круг. В ней вычисленная скорость второго автомобиля $66$ км/ч действительно меньше скорости первого $84$ км/ч. Это противоречит условию задачи.

\subtopic{Ключевая идея}
\textcolor{blue!60!black}{\textbf{Ключевая идея:}} при движении по кругу в одном направлении разность пройденных путей при первом обгоне на круг равна длине трассы:
\[
 (v_{\text{быстр.}}-v_{\text{медл.}})t=L.
\]
Разность скоростей берётся в порядке «более быстрый минус более медленный».

\subtopic{Исправление от последнего правильного шага}
Сохраняем уже найденные Вами расстояния $56$ км и $\dfrac{2x}{3}$ км. Меняем порядок вычитания:
\begin{align*}
 \frac{2x}{3}-56&=12,\\
 \frac{2x}{3}&=12+56=68,\\
 2x&=204,\\
 x&=102.
\end{align*}

\subtopic{Полное правильное решение}
\textbf{Шаг 1.} Обозначим скорость второго автомобиля через $x$ км/ч. Поскольку обгоняет второй автомобиль, $x>84$.

\textbf{Шаг 2.} Переведём время в часы:
\[
 t=40\text{ мин}=\frac{40}{60}=\frac23\text{ ч}.
\]
\textbf{Шаг 3.} За это время первый автомобиль проехал:
\[
 S_1=84\cdot\frac23=56\text{ км}.
\]
Второй автомобиль проехал $S_2=\dfrac{2x}{3}$ км.

\textbf{Шаг 4.} Так как второй автомобиль впервые обогнал первый на один круг, разница их путей составляет $12$ км:
\[
 \frac{2x}{3}-56=12.
\]
\textbf{Шаг 5.} Решим уравнение:
\[
 \frac{2x}{3}=68 \quad\Longrightarrow\quad
 x=68\cdot\frac32=102\text{ км/ч}.
\]
\clearpage
\subtopic{Дополнительный способ через относительную скорость}
За $\dfrac23$ часа второй автомобиль должен \emph{набрать} преимущество в $12$ км. Поэтому разность скоростей равна
\[
 \Delta v=\frac{12}{2/3}=18\text{ км/ч},\qquad
 x=84+18=102\text{ км/ч}.
\]

\subtopic{Проверка}
За $40$ минут первый автомобиль проходит $84\cdot\dfrac23=56$ км, второй --- $102\cdot\dfrac23=68$ км. Разность путей:
\[
 68-56=12\text{ км}.
\]
Она равна ровно одному кругу. Поскольку разность скоростей постоянна и положительна, до $40$-й минуты разность пройденных путей была меньше $12$ км. Значит, это действительно \emph{первый} обгон на полный круг.

\subtopic{Итог}
\textcolor{teal!60!black}{\textbf{Итог:}} скорость второго автомобиля составляет $\boxed{102\text{ км/ч}}$. Главный навык --- определить по тексту задачи, чей путь должен быть больше, и лишь затем записывать разность.

\subtopic{Задача на закрепление}
Длина круговой трассы составляет $15$ км. Из одной точки одновременно стартовали два автомобиля, движущиеся в одном направлении. Скорость первого автомобиля равна $60$ км/ч, скорость второго постоянна. Через $30$ минут второй автомобиль впервые обогнал первый ровно на один полный круг. Найдите скорость второго автомобиля. Ответ дайте в километрах в час.

\vspace{13mm}
\begin{center}
{\small\color{textgray} Лёвин Артём Александрович эксклюзивно для Анны Банновой}
\end{center}
\end{document}